\documentclass[UTF8,zihao=-4]{ctexart}
\usepackage[a4paper,margin=2.2cm]{geometry}
\usepackage{xcolor}
\usepackage{hyperref}
\usepackage{enumitem}
\usepackage{fontspec}
\usepackage{amsmath}
\IfFontExistsTF{Microsoft YaHei}{\setCJKmainfont{Microsoft YaHei}}{\setCJKmainfont{SimSun}[BoldFont=SimHei]}
\IfFontExistsTF{Microsoft YaHei UI}{\setCJKsansfont{Microsoft YaHei UI}}{}
\IfFontExistsTF{Consolas}{\setmonofont{Consolas}}{}
\definecolor{linkblue}{RGB}{30,80,145}
\hypersetup{colorlinks=true,linkcolor=linkblue,urlcolor=linkblue,pdfborder={0 0 0}}
\setlength{\parindent}{2em}
\setlength{\parskip}{0.25em}
\linespread{1.16}
\setlist{itemsep=0.2em,topsep=0.35em,leftmargin=2.5em}
\pagestyle{plain}

\title{12\_动态规划}

\author{}

\date{}

\begin{document}

\maketitle

\section*{12 动态规划}

\subsection*{题目原文}

动态规划的思想和特点，针对经典问题能写出子问题的递归式，描述算法思想和伪代码：最长公共子序列相关问题、抽牌拿取最大值问题、回文判断/分割相关问题。

\subsection*{考查类型}

概念题和算法设计题。

\subsection*{动态规划思想}

动态规划用于具有以下性质的问题：

\begin{quote}
\small\ttfamily\raggedright
\relax 最优子结构\\
\relax 重叠子问题\\
\end{quote}

做法是先定义状态，再写状态转移，最后确定计算顺序。

答题四步：

\begin{enumerate}
\item 定义 \texttt{dp} 状态含义。
\item 写初始条件。
\item 写状态转移方程。
\item 说明遍历顺序和最终答案。
\end{enumerate}

\subsection*{最长公共子序列}

给定字符串 \texttt{X[1..m]} 和 \texttt{Y[1..n]}，令：

\begin{quote}
\small\ttfamily\raggedright
\relax dp[i][j]\ =\ X[1..i]\ 和\ Y[1..j]\ 的最长公共子序列长度\\
\end{quote}

初始条件：

\begin{quote}
\small\ttfamily\raggedright
\relax dp[0][j]\ =\ 0\\
\relax dp[i][0]\ =\ 0\\
\end{quote}

状态转移：

\begin{quote}
\small\ttfamily\raggedright
\relax if\ X[i]\ ==\ Y[j]:\\
\relax \ \ \ \ dp[i][j]\ =\ dp[i-1][j-1]\ +\ 1\\
\relax else:\\
\relax \ \ \ \ dp[i][j]\ =\ max(dp[i-1][j],\ dp[i][j-1])\\
\end{quote}

答案：

\begin{quote}
\small\ttfamily\raggedright
\relax dp[m][n]\\
\end{quote}

时间复杂度：

\begin{quote}
\small\ttfamily\raggedright
\relax O(mn)\\
\end{quote}

空间复杂度：

\begin{quote}
\small\ttfamily\raggedright
\relax O(mn)\\
\end{quote}

若只要求长度，可以滚动数组优化为 $O(n)$ 空间。

\subsection*{抽牌拿取最大值问题}

PDF 未给出完整题面。以下小节只分析本文件采用的练习模型：一排牌，每次只能从左端或右端取一张，两人都采用最优策略，求先手能取得的最大收益或最大分差。课堂题面不同，则需要按课堂题面重写该小节。

定义：

\begin{quote}
\small\ttfamily\raggedright
\relax dp[i][j]\ =\ 当前行动者面对区间\ cards[i..j]\ 时，相对对手最多领先多少分\\
\end{quote}

状态转移：

\begin{quote}
\small\ttfamily\raggedright
\relax dp[i][j]\ =\ max(\\
\relax \ \ \ \ cards[i]\ -\ dp[i+1][j],\\
\relax \ \ \ \ cards[j]\ -\ dp[i][j-1]\\
\relax )\\
\end{quote}

含义：

当前行动者若取左端 \texttt{cards[i]}，接下来对手面对 \texttt{[i+1..j]}，对手相对当前行动者的领先为 \texttt{dp[i+1][j]}，所以当前行动者净领先为 \texttt{cards[i] - dp[i+1][j]}。

初始条件：

\begin{quote}
\small\ttfamily\raggedright
\relax dp[i][i]\ =\ cards[i]\\
\end{quote}

计算顺序：

\begin{quote}
\small\ttfamily\raggedright
\relax 区间长度从\ 1\ 到\ n\\
\end{quote}

答案：

\begin{quote}
\small\ttfamily\raggedright
\relax dp[0][n-1]\\
\end{quote}

若需要先手最大总分，可结合总分 $S$ 得到：

\begin{quote}
\small\ttfamily\raggedright
\relax first\ =\ (S\ +\ dp[0][n-1])\ /\ 2\\
\end{quote}

\subsection*{回文判断}

给定字符串 \texttt{s[0..n-1]}，令：

\begin{quote}
\small\ttfamily\raggedright
\relax pal[i][j]\ =\ s[i..j]\ 是否为回文串\\
\end{quote}

初始和转移：

\begin{quote}
\small\ttfamily\raggedright
\relax pal[i][i]\ =\ true\\
\relax pal[i][i+1]\ =\ (s[i]\ ==\ s[i+1])\\
\mbox{}\\
\relax if\ s[i]\ ==\ s[j]\ and\ pal[i+1][j-1]:\\
\relax \ \ \ \ pal[i][j]\ =\ true\\
\relax else:\\
\relax \ \ \ \ pal[i][j]\ =\ false\\
\end{quote}

计算顺序：

\begin{quote}
\small\ttfamily\raggedright
\relax 按区间长度从短到长\\
\end{quote}

\subsection*{回文分割}

目标：将字符串切成若干回文子串，求最少切割次数。

先用上面的 \texttt{pal[i][j]} 预处理任意区间是否回文。

令：

\begin{quote}
\small\ttfamily\raggedright
\relax cut[i]\ =\ s[0..i]\ 的最少切割次数\\
\end{quote}

转移：

\begin{quote}
\small\ttfamily\raggedright
\relax if\ pal[0][i]:\\
\relax \ \ \ \ cut[i]\ =\ 0\\
\relax else:\\
\relax \ \ \ \ cut[i]\ =\ min(cut[j]\ +\ 1)\ for\ 0\ ≤\ j\ <\ i\ and\ pal[j+1][i]\\
\end{quote}

时间复杂度：

\[O(n^{2})\]

\subsection*{例题}

\subsubsection*{例题 1}

题目：求字符串 $X = ABCBDAB$ 和 $Y = BDCABA$ 的最长公共子序列长度。

解答：

令：

\begin{quote}
\small\ttfamily\raggedright
\relax dp[i][j]\ =\ X[1..i]\ 和\ Y[1..j]\ 的最长公共子序列长度\\
\end{quote}

按转移方程填表后得到：

\begin{quote}
\small\ttfamily\raggedright
\relax dp[7][6]\ =\ 4\\
\end{quote}

因此最长公共子序列长度为 \texttt{4}。一个长度为 4 的公共子序列为：

\begin{quote}
\small\ttfamily\raggedright
\relax BCBA\\
\end{quote}

\subsubsection*{例题 2}

题目：一排牌为：

\begin{quote}
\small\ttfamily\raggedright
\relax [3,\ 9,\ 1,\ 2]\\
\end{quote}

每次只能从左端或右端取牌，两人都采用最优策略。求先手最终得分。

解答：

沿用上文定义：

\begin{quote}
\small\ttfamily\raggedright
\relax dp[i][j]\ =\ 当前行动者面对\ cards[i..j]\ 时，相对对手最多领先多少分\\
\end{quote}

单元素区间：

\begin{quote}
\small\ttfamily\raggedright
\relax dp[0][0]\ =\ 3\\
\relax dp[1][1]\ =\ 9\\
\relax dp[2][2]\ =\ 1\\
\relax dp[3][3]\ =\ 2\\
\end{quote}

长度为 2：

\begin{quote}
\small\ttfamily\raggedright
\relax dp[0][1]\ =\ max(3\ -\ 9,\ 9\ -\ 3)\ =\ 6\\
\relax dp[1][2]\ =\ max(9\ -\ 1,\ 1\ -\ 9)\ =\ 8\\
\relax dp[2][3]\ =\ max(1\ -\ 2,\ 2\ -\ 1)\ =\ 1\\
\end{quote}

长度为 3：

\begin{quote}
\small\ttfamily\raggedright
\relax dp[0][2]\ =\ max(3\ -\ 8,\ 1\ -\ 6)\ =\ -5\\
\relax dp[1][3]\ =\ max(9\ -\ 1,\ 2\ -\ 8)\ =\ 8\\
\end{quote}

长度为 4：

\begin{quote}
\small\ttfamily\raggedright
\relax dp[0][3]\ =\ max(3\ -\ 8,\ 2\ -\ (-5))\ =\ 7\\
\end{quote}

总分为：

\[3 + 9 + 1 + 2 = 15\]

先手得分为：

\[(15 + 7) / 2 = 11\]

\subsubsection*{例题 3}

题目：对字符串 \texttt{aab} 求最少回文分割次数。

解答：

回文子串包括：

\begin{quote}
\small\ttfamily\raggedright
\relax a\\
\relax a\\
\relax b\\
\relax aa\\
\end{quote}

最优分割为：

\begin{quote}
\small\ttfamily\raggedright
\relax aa\ |\ b\\
\end{quote}

切割次数为：

\begin{quote}
\small\ttfamily\raggedright
\relax 1\\
\end{quote}

用 \texttt{cut} 表示：

\begin{quote}
\small\ttfamily\raggedright
\relax cut[0]\ =\ 0\ \ \ \ \ \ \ \ a\\
\relax cut[1]\ =\ 0\ \ \ \ \ \ \ \ aa\\
\relax cut[2]\ =\ cut[1]\ +\ 1\ =\ 1\\
\end{quote}

所以答案为 \texttt{1}。

\subsection*{常见误区}

动态规划不是只写递归式。完整答案还要说明状态含义、边界条件、计算顺序和最终答案。

区间动态规划要按区间长度递增计算，避免使用尚未计算出的状态。

最长公共子序列不是最长公共子串。子序列允许不连续，子串必须连续。

\end{document}
